% 元の
% \documentclass[lualatex,book,paper=a4paper,jafontsize=11.5pt,jafontscale=1.05]{jlreq}
\documentclass[lualatex,paper=a4paper,jafontsize=11.5pt,jafontscale=1.05]{jlreq}
% 2ページ目や章の前に一ページ余分なものがあるのが気になる場合は openanyを\documentclass[...,book,openany ,..]として追加する

%フォント:
\usepackage{luatexja-fontspec}
\usepackage{hack-fonts}

%表紙用:
\usepackage{hack-cover}
\usepackage{enumitem}
% 画像
\usepackage{graphicx}
\usepackage{float}
\usepackage{array}
\usepackage{makecell}

% リンク
\usepackage[hidelinks]{hyperref}

% ソースコード
\usepackage{listings}
\lstset{
  frame=single,
  basicstyle=\small\ttfamily,
  tabsize=4,
  breaklines=true
}

% 色・ガントチャート
\usepackage{xcolor}
\usepackage{pgfgantt}

\begin{document}
\section*{作業ログ}
\noindent
報告者:恩田 隼士\\
報告日: \today \\
プロジェクト名: 単一のフルカラーLEDとArduinoを用いた光無線通信による楽器間同期演奏システム \\
\hrulefill

\section{進捗サマリー（概要）}
\begin{itemize}
 \item 全体進捗率: 85\%
 \item マイルストーン達成状況: 一部達成
 \item 主要成果:
 \begin{itemize}
    \item 演奏者Arduino4台と指揮者Arduino1台を用いた実機演奏動作の確認
    \item 演奏停止プログラム（点滅未検知タイムアウト機能）の実装
    \item アンケート用ヴァイオリン音源の録音およびアンケート内容のチーム内すり合わせ
 \end{itemize}
 \item 重大課題（影響度）: なし（残課題は\ref{sec:issues}節：課題管理を参照）
\end{itemize}

\noindent
※今週新たに「遅延テスト」「4台同時演奏テスト」「プレゼン準備」を作業項目として追加したため，
分母の増加により全体進捗率は前回（95\%）から見かけ上低下している．

\section{作業ログ（詳細トラッキング）}
\begin{table}[H]
  \centering
  \caption{作業ログ}
   \scalebox{0.7}{
    \begin{tabular}{lcccccc} \hline
      作業項目 & 開始日 & 予定完了日 & 状態 & 進捗率 & 依存関係・ブロッカー & コメント \\ \hline
      楽器演奏 & 05/20 & 05/23 & 完了 & 100\% & なし & - \\
      演奏開始・停止ボタン（設計） & 05/23 & 05/27 & 完了 & 100\% & 部品入手 & - \\
      楽器音 & 05/23 & 05/27 & 完了 & 100\% & 楽器演奏 & - \\
      楽譜配列 & 05/26 & 05/28 & 完了 & 100\% & 楽器音 & - \\
      ヴァイオリン音 & 05/28 & 06/12 & 完了 & 100\% & 楽器音 & - \\
      演奏開始・停止ボタン（実装） & 05/28 & 06/09 & 完了 & 100\% & 設計 & - \\
      結合 & 06/17 & 07/07 & 進行中 & 90\% & 単体実装 & \makecell[l]{4台演奏者Arduino＋指揮者1台での\\実機演奏動作を確認．\\安定性向上は継続中} \\
      開始停止スイッチテスト & 06/13 & 06/19 & 完了 & 100\% & 実装の完了 & - \\
      楽器音テスト & 06/13 & 06/19 & 完了 & 100\% & 実装の完了 & - \\
      輪唱演奏テスト & 06/17 & 06/25 & 進行中 & 60\% & 結合の完了 & - \\
      演奏停止プログラム & 06/29 & 06/29 & 完了 & 100\% & 結合 & \makecell[l]{点滅未検知タイムアウトによる\\演奏状態リセット機能を実装} \\
      遅延テスト & 06/30 & 07/07 & 未着手 & 0\% & 結合の完了 & 新規追加 \\
      4台同時演奏テスト & 06/30 & 07/07 & 進行中 & 30\% & 結合の完了 & \makecell[l]{動作確認は実施．\\安定動作の検証はこれから} \\
      プレゼン準備 & 06/30 & 07/07 & 未着手 & 0\% & テストの完了 & 新規追加 \\
      \hline
    \end{tabular}
  }
\end{table}


% ===== 計画前ガントチャート =====
\subsection{担当箇所のガントチャート（計画前）}
\noindent
\begin{figure}[H]
\centering
\resizebox{0.92\linewidth}{!}{%
\begin{ganttchart}[
  hgrid,
  vgrid=dotted,
  time slot format=isodate,
  time slot unit=day,
  x unit=0.52cm,
  y unit title=0.6cm,
  y unit chart=0.65cm,
  title label font=\tiny,
  bar label font=\small,
  bar label node/.append style={text width=4.8cm, align=right},
  bar height=0.55,
  bar top shift=0.23,
  bar/.append style={fill=violet!45, draw=violet!75},
]{2026-05-20}{2026-06-16}
  \gantttitlecalendar{month=shortname} \\
  \gantttitlecalendar{day} \\
  \ganttbar{楽器演奏}{2026-05-20}{2026-05-23} \\
  \ganttbar{演奏開始・停止ボタン}{2026-05-23}{2026-05-27} \\
  \ganttbar{楽器音}{2026-05-23}{2026-05-27} \\
  \ganttbar{楽譜配列}{2026-05-26}{2026-05-28} \\
  \ganttbar{ヴァイオリン音}{2026-05-28}{2026-05-31}
  \ganttbar[bar label node/.append style={opacity=0}]{}{2026-06-04}{2026-06-09} \\
  \ganttbar{開始停止スイッチテスト}{2026-06-10}{2026-06-16} \\
  \ganttbar{楽器音テスト}{2026-06-10}{2026-06-16}
\end{ganttchart}%
}
\caption{担当箇所ガントチャート（計画前）}
\end{figure}

% ===== 現在のガントチャート =====
% 凡例：■ 灰色 = 完了，■ 紫色 = 未完了 / 進行中
% タスクの状態に応じて fill=gray!60（完了）または fill=violet!50（未完了）を変更
\subsection{担当箇所のガントチャート（現在・進捗状況）}
\noindent
{\small \textcolor{gray!50}{\rule{1em}{0.8em}}\ 完了\quad
       \textcolor{violet!60}{\rule{1em}{0.8em}}\ 未完了／進行中}
\begin{figure}[H]
\centering
\resizebox{0.92\linewidth}{!}{%
\begin{ganttchart}[
  hgrid,
  vgrid=dotted,
  time slot format=isodate,
  time slot unit=day,
  x unit=0.52cm,
  y unit title=0.6cm,
  y unit chart=0.65cm,
  title label font=\tiny,
  bar label font=\small,
  bar label node/.append style={text width=4.8cm, align=right},
  bar height=0.55,
  bar top shift=0.23,
]{2026-05-20}{2026-07-07}
  \gantttitlecalendar{month=shortname} \\
  \gantttitlecalendar{day} \\
  % ---- 完了タスク (gray!60) ----
  \ganttbar[bar/.append style={fill=gray!60, draw=gray!80}]
    {楽器演奏}{2026-05-20}{2026-05-23} \\
  \ganttbar[bar/.append style={fill=gray!60, draw=gray!80}]
    {演奏開始・停止ボタン（設計）}{2026-05-23}{2026-05-27} \\
  \ganttbar[bar/.append style={fill=gray!60, draw=gray!80}]
    {楽器音}{2026-05-23}{2026-05-27} \\
  \ganttbar[bar/.append style={fill=gray!60, draw=gray!80}]
    {楽譜配列}{2026-05-26}{2026-05-28} \\
  % ---- ヴァイオリン音：完了 (gray!60) ----
  \ganttbar[bar/.append style={fill=gray!60, draw=gray!80}]
    {ヴァイオリン音}{2026-05-28}{2026-05-31}
  \ganttbar[bar/.append style={fill=gray!60, draw=gray!80},
    bar label node/.append style={opacity=0}]
    {}{2026-06-04}{2026-06-12} \\
  \ganttbar[bar/.append style={fill=gray!60, draw=gray!80}]
    {演奏開始・停止ボタン（実装）}{2026-05-28}{2026-05-31}
  \ganttbar[bar/.append style={fill=gray!60, draw=gray!80},
    bar label node/.append style={opacity=0}]
    {}{2026-06-04}{2026-06-09} \\
  \ganttbar[bar/.append style={fill=gray!60, draw=gray!80}]
    {開始停止スイッチテスト}{2026-06-13}{2026-06-19} \\
  \ganttbar[bar/.append style={fill=gray!60, draw=gray!80}]
    {楽器音テスト}{2026-06-13}{2026-06-19} \\
  % ---- 輪唱演奏（実機結合）：継続 ----
  \ganttbar[bar/.append style={fill=violet!50, draw=violet!80}]
    {輪唱演奏（実機結合）}{2026-06-17}{2026-07-07} \\
  % ---- 今週完了：演奏停止プログラム ----
  \ganttbar[bar/.append style={fill=gray!60, draw=gray!80}]
    {演奏停止プログラム}{2026-06-29}{2026-06-29} \\
  % ---- 今週追加：遅延テスト ----
  \ganttbar[bar/.append style={fill=violet!50, draw=violet!80}]
    {遅延テスト}{2026-06-30}{2026-07-07} \\
  % ---- 今週追加：4台同時演奏テスト ----
  \ganttbar[bar/.append style={fill=violet!50, draw=violet!80}]
    {4台同時演奏テスト}{2026-06-30}{2026-07-07} \\
  % ---- 今週追加：プレゼン準備 ----
  \ganttbar[bar/.append style={fill=violet!50, draw=violet!80}]
    {プレゼン準備}{2026-06-30}{2026-07-07}
\end{ganttchart}%
}
\caption{担当箇所ガントチャート（現在：灰＝完了，紫＝未完了）}
\end{figure}

ガントチャートにおいて，06/01$\sim$06/03の期間は中間試験対策期間として空白にしている．
なお，計画前ガントチャートは計画当初の予定を保持するため，以降の計画変更は反映していない．現在の進捗状況との差分は，「現在・進捗状況」ガントチャートとの比較により確認できる．

% ===== 日ごとの作業時間・内容 =====
\subsection{日ごとの作業時間・内容}
\begin{table}[H]
  \centering
  \caption{日ごとの作業時間・内容}
  \scalebox{0.88}{
    \begin{tabular}{ccp{7.5cm}} \hline
      日付 & 作業時間 & 作業内容・備考 \\ \hline
      2026/06/24 & 4h & 演奏者Arduino4台＋指揮者Arduino1台による実機演奏動作確認．LED閾値が一定に定まらない問題および演奏停止プログラム未実装の課題を発見したため，詳細テストには未到達 \\[4pt]
      2026/06/27 & 1h & アンケート用ヴァイオリン音源の録音 \\[4pt]
      2026/06/29 & 2h & チーム内アンケート内容すり合わせ・演奏停止プログラム（点滅未検知タイムアウト機能）の作成 \\

        % ※必要に応じて行を追加
      \hline
    \end{tabular}
  }

\end{table}

% ===== 今週追加：ログの詳細（TPM・TRL） =====
\section{ログの詳細}

\subsection{演奏者Arduino4台＋指揮者Arduino1台による実機演奏動作確認}
前回までに演奏者Arduino2台での輪唱演奏動作を確認していたため，
今週は本来の構成である演奏者Arduino4台と指揮者Arduino1台を接続した実機演奏動作確認を実施した．
4声輪唱が成立すること自体は確認できたが，以下2点の課題により詳細テスト（伝送遅延計測等）には到達していない．
\begin{itemize}
  \item LEDの色判定閾値が環境（USB電源・周囲光）により一定に定まらず，4台同時起動時に検出が不安定となる場面が散見された．閾値設計の見直しは別担当が継続対応中である．
  \item 楽譜末尾以外のタイミング（指揮者切断・配線外れ等）で演奏状態を待機に戻す仕組みが存在せず，再演奏の度に演奏者Arduinoの再起動が必要であった．
\end{itemize}

\subsection{演奏停止プログラム（点滅未検知タイムアウト機能）}
上記2点目の課題に対応するため，演奏者Arduinoに「点滅未検知タイムアウト」機能を追加実装した．
指揮者LEDの立ち上がりエッジが一定時間検出されない場合，演奏進行状態（楽譜進行カウンタ noteNum，歌詞インデックス lyricIndex，色判定ラッチ armed）を自動的に待機状態へリセットする．
タイムアウト値は3000\,msとし，BPM30（最遅想定）における半拍周期1000\,msの3倍を確保した．
検出系状態（inBlink・clearMax・blinkTime1・blinkInterval）は次の立ち上がりエッジで自然に再初期化されるためリセット対象から除外し，副作用を最小化している．
また，起動直後の誤発火を避けるため，発火条件には「ラッチ済み（armed=true）または1音以上演奏済み（noteNum$>$0）」のガードを設けている．
本機能は \texttt{feature/player-blink-timeout} ブランチで実装し，\texttt{playerRGBMethod} 配下の実装に統合済みである．

\subsection{アンケート関連作業}
プレゼンに向けた被験者アンケートの準備として，ヴァイオリン音源の録音（06/27）と
チーム内でのアンケート内容すり合わせ（06/29）を実施した．
アンケート設計・実施はプレゼン準備項目に内包する形で次週以降進行する．

\subsection{技術性能指標（TPM）}
TPMはシステムの最終的な性能に関わる技術を管理するための指標である．
本プロジェクトの各指標のうち，自分の担当は「音色の再現性」である．
音色の再現性については前回報告にて全4特徴量（AT，SC，DSV，OER）が目標値を達成済みである．
今週は4台実機演奏確認・演奏停止プログラム実装フェーズであり，新たなTPM計測は実施していない．
次回以降，結合後の伝送遅延計測等を実施予定である．

\subsection{技術成熟度（TRL）}
TRLは機能ごとの進捗を確認する指標である．定義を表\ref{tab:trl}に示す．

\begin{table}[H]
  \centering
  \caption{TRLの定義}
  \label{tab:trl}
  \scalebox{0.9}{
    \begin{tabular}{cp{10cm}} \hline
      TRL & 定義 \\ \hline
      1 & 概念レベル：アイディアや原理が確立されており，説明できるレベル \\
      2 & 基本動作レベル：単体で動作確認が可能なレベル \\
      3 & 結合動作レベル：機能を結合し，実際の使用環境に近い環境で動作確認が可能なレベル \\
      4 & 実運用レベル：実際の使用環境で安定的な稼働が確認できるレベル \\
      \hline
    \end{tabular}
  }
\end{table}

自分の担当機能の現状を表\ref{tab:trl_status}に示す．
4台＋指揮者の実機構成で演奏動作自体は確認できたものの，LED閾値の不安定さに起因して安定動作には至っていないため，
同期演奏システムの結合および輪唱演奏はいずれもTRL3継続と判定した．
演奏停止プログラムは単体動作の確認が完了しており，TRL2と判定した．

\begin{table}[H]
  \centering
  \caption{担当機能のTRL現状}
  \label{tab:trl_status}
  \scalebox{0.9}{
    \begin{tabular}{lcp{6cm}} \hline
      担当機能 & 現在のTRL & 根拠 \\ \hline
      同期演奏システムの結合 & TRL3 & 4台＋指揮者での実機演奏動作を確認．安定性向上は継続中 \\
      輪唱演奏 & TRL3 & 4台での輪唱演奏動作を確認．安定性向上は継続中 \\
      演奏停止プログラム & TRL2 & 単体での動作確認まで完了 \\
      \hline
    \end{tabular}
  }
\end{table}

\section{課題管理}\label{sec:issues}

\subsection{設計書要件との適合確認}
\begin{itemize}
  \item 原因: 結合・輪唱実装の過程で，設計書との整合確認が一部未完了であった．
  \item 影響度: 中
  \item 優先度: 中
  \item 対応策: チーム内で実装仕様と設計書記載を突き合わせ，追記・修正箇所を確定した．
  \item 期限: 2026/06/30
  \item 状態: 完了（設計書への反映済み）
\end{itemize}

\subsection{LED色判定閾値の環境依存}
\begin{itemize}
  \item 原因: TCS34725による色判定の閾値が，USB電源の変動および周囲光環境により一定に定まらず，4台同時起動時に検出が不安定となる．
  \item 影響度: 中
  \item 優先度: 高
  \item 対応策: 別担当が閾値設計の見直し（演奏比率計算方式・チューニング機能等）を継続実施中．本ログ報告者側は実装結果を受け取り次第，結合テストへ反映する．
  \item 期限: 2026/07/07
  \item 状態: 進行中（別担当対応中）
\end{itemize}

\section{AI利用状況と検証}
\subsection{利用内容}
\subsubsection{演奏者Arduinoの演奏停止プログラム（点滅未検知タイムアウト機能）}
\begin{itemize}
    \item 利用目的：指揮者切断・配線外れ等で点滅信号が途絶えた場合に，演奏者Arduinoの演奏進行状態を自動的に待機状態へ戻し，再演奏時の再起動操作を不要とする機能の実装
    \item 入力（プロンプト要旨）：
    \begin{itemize}
        \item 既存の \texttt{playerRGBMethod.ino} に対し，指揮者LEDの立ち上がりエッジが一定時間検出されない場合に演奏状態を最初に戻す処理を追加したい旨を指定した．
        \item リセット対象は楽譜進行状態（noteNum，lyricIndex，armed）のみとし，検出系状態およびWeb経由の設定値（楽器・音量・楽譜ID）は保持することを条件として与えた．
        \item BPM下限値および半拍周期から逆算した妥当なタイムアウト値の根拠も合わせて提示するよう依頼した．
        \item 起動直後（1音も鳴っていない状態）での誤発火を避けるガード条件を設けることを指定した．
    \end{itemize}
    \item 出力概要：
    \begin{itemize}
        \item タイムアウト定数 \texttt{BLINK\_TIMEOUT\_MS = 3000} を導入し，BPM30の半拍周期1000\,msに対して3倍の余裕を持たせる設計根拠が得られた．
        \item メインループ内で \texttt{blinkTime1 != 0 \&\& (armed || noteNum > 0)} のガード条件下で経過時間判定を行い，超過時に楽譜進行状態のみをリセットする実装が得られた．
        \item 検出系状態は次の立ち上がりエッジで自然初期化されるため，リセット対象から除外する設計判断が得られた．
        \item タイムアウト直後の再開時に最初の音の \texttt{blinkInterval} が一時的に長くなる副作用が許容範囲である旨も合わせて確認できた．
    \end{itemize}
\end{itemize}

\subsection{検証方法}
% ※第三者が再現できるレベルで記述
\begin{itemize}
  \item 検証方法:
  \begin{itemize}
      \item 実機（演奏者Arduino・指揮者Arduino・Processing PC）での動作確認による．
  \end{itemize}
  \item 参照資料:
  \begin{itemize}
      \item 本プロジェクト設計書\footnote{チーム内部資料．}．
      \item Arduino UNO R4 WiFi 公式リファレンス（WiFiS3 ライブラリ）\cite{wifis3}．
  \end{itemize}
  \item 検証手順：
  \begin{enumerate}
      \item 演奏者Arduinoを起動し，指揮者LEDの自パート色を1回以上受光させて演奏状態に遷移させる．
      \item 指揮者Arduinoの電源を切断し，演奏者Arduino側の演奏進行が \texttt{BLINK\_TIMEOUT\_MS}（3000\,ms）経過後に待機状態へ戻ることを確認する．
      \item 再度指揮者Arduinoを起動し，演奏者Arduinoの再起動なしで自パート色受光から演奏が再開できることを確認する．
      \item 起動直後（1音も鳴らしていない状態）で3000\,ms以上待機した際に，誤ってリセット処理が発火しないことを確認する．
  \end{enumerate}
\end{itemize}

\subsection{ファクトチェック}
\begin{itemize}
  \item 一致点：
  \begin{itemize}
    \item 本プロジェクトのBPM想定範囲（30〜240）における半拍周期に対し，3秒のタイムアウト値は十分な余裕を持つ．
    \item リセット対象を楽譜進行状態（noteNum，lyricIndex，armed）に限定する方針は，既存の色判定ラッチ仕様（[[project-color-latch-spec]]）と整合する．
    \item Web経由の設定値（楽器・音量・楽譜ID）をリセット対象から除外する判断は，ユーザ選択保持の観点から妥当である．
  \end{itemize}
  \item 不一致・疑義：なし
  \item 最終判断：採用
  \item 根拠：
  \begin{itemize}
    \item 修正後の \texttt{playerRGBMethod.ino} を実機投入し，指揮者切断時に演奏者Arduinoが自動的に待機状態へ戻ること，および再起動なしで演奏再開が可能となることを確認した．
    \item 起動直後の待機状態において誤発火が発生しないことを確認した．
  \end{itemize}
\end{itemize}

\section{次回計画（次回までに実施する内容）}
\begin{itemize}
  \item 実施事項：
  \begin{itemize}
    \item 4台同時演奏テスト：LED閾値設計の改善版を反映した上で，演奏者Arduino4台での輪唱演奏が安定動作することを確認する（TRL4を目指す）．
    \item 遅延テスト：指揮者LEDの点滅から演奏者Arduino受光・Processing発音までの伝送遅延を計測し，許容範囲内であることを確認する．
    \item プレゼン準備：被験者アンケートの設計・実施およびプレゼン資料の作成を進める．
  \end{itemize}
  \item 達成基準：
  \begin{itemize}
    \item 4台すべてで輪唱が安定して成立すること．
    \item 遅延が演奏上知覚されない範囲に収まること．
    \item プレゼン資料・アンケート結果の整備が完了していること．
  \end{itemize}
  \item 期限：2026/07/07
\end{itemize}

\section{その他}
連絡事項・懸念点：なし

\section{付録}
添付資料を以下のリポジトリURLに示す．添付範囲は全体である．\\
\url{https://github.com/crab2424/LEDplaying_system}



\begin{thebibliography}{9}
\bibitem{processing} Processing Foundation (2026), "Processing Reference". \url{https://processing.org/reference/}. (2026年6月30日 参照)
\bibitem{minim_unpatch} D. Di Fede (2026), "Minim : UGen.unpatch()". \url{http://code.compartmental.net/minim/ugen_method_unpatch.html}. (2026年6月30日 参照)
\bibitem{wifis3} Arduino (2026), "WiFiS3 Library". \url{https://docs.arduino.cc/libraries/wifi/}. (2026年6月30日 参照)
\end{thebibliography}

\end{document}
